% ============================================================================
% CHAPTER 8: SECURITY FEATURES & ROADMAP
% ============================================================================

\chapter{Security Features \& Roadmap}

\section{Current Security Implementation}

\subsection{Authentication \& Authorization}

\subsubsection{JWT Token-Based Authentication}

The platform implements secure, stateless authentication using JSON Web Tokens (JWT).

\textbf{Token Architecture:}
\begin{lstlisting}[style=pythonstyle]
# services/api-gateway/middleware/auth.py
from jose import JWTError, jwt
from datetime import datetime, timedelta

class JWTAuthentication:
    def __init__(self):
        self.secret_key = os.getenv("JWT_SECRET_KEY")
        self.algorithm = "HS256"
        self.access_token_expire = 1800  # 30 minutes
        self.refresh_token_expire = 604800  # 7 days
    
    def create_access_token(self, user_id: str, role: str) -> str:
        """Generate short-lived access token"""
        payload = {
            "sub": user_id,
            "role": role,
            "type": "access",
            "exp": datetime.utcnow() + timedelta(seconds=self.access_token_expire),
            "iat": datetime.utcnow()
        }
        return jwt.encode(payload, self.secret_key, algorithm=self.algorithm)
    
    def create_refresh_token(self, user_id: str) -> str:
        """Generate long-lived refresh token"""
        payload = {
            "sub": user_id,
            "type": "refresh",
            "exp": datetime.utcnow() + timedelta(seconds=self.refresh_token_expire),
            "iat": datetime.utcnow()
        }
        return jwt.encode(payload, self.secret_key, algorithm=self.algorithm)
    
    def verify_token(self, token: str) -> dict:
        """Verify and decode JWT token"""
        try:
            payload = jwt.decode(token, self.secret_key, algorithms=[self.algorithm])
            return payload
        except JWTError:
            raise HTTPException(status_code=401, detail="Invalid token")
\end{lstlisting}

\textbf{Token Security Features:}
\begin{itemize}
    \item \textbf{Short Expiry}: Access tokens expire in 30 minutes
    \item \textbf{Refresh Mechanism}: Refresh tokens valid for 7 days
    \item \textbf{Secure Storage}: HTTP-only cookies (not accessible to JavaScript)
    \item \textbf{CSRF Protection}: SameSite cookie attribute set to 'Lax'
    \item \textbf{Token Blacklisting}: Revoked tokens stored in Redis
\end{itemize}

\subsubsection{Role-Based Access Control (RBAC)}

Comprehensive permission system with 5 predefined roles.

\begin{table}[h]
\centering
\caption{RBAC Permission Matrix}
\begin{tabularx}{\textwidth}{lXXXXX}
\toprule
\textbf{Resource} & \textbf{Admin} & \textbf{Data Mgr} & \textbf{Reviewer} & \textbf{Analyst} & \textbf{Guest} \\
\midrule
Upload Documents & ✓ & ✓ & ✗ & ✗ & ✗ \\
View Documents & ✓ & ✓ & ✓ & ✓ & ✓ \\
Classify Documents & ✓ & ✓ & ✗ & ✗ & ✗ \\
Review ICSRs & ✓ & ✓ & ✓ & ✗ & ✗ \\
Approve ICSRs & ✓ & ✗ & ✓ & ✗ & ✗ \\
Manage Users & ✓ & ✗ & ✗ & ✗ & ✗ \\
View Knowledge Hub & ✓ & ✓ & ✓ & ✓ & ✓ \\
Investigation Boards & ✓ & ✓ & ✓ & ✓ (read) & ✗ \\
System Settings & ✓ & ✗ & ✗ & ✗ & ✗ \\
API Access & ✓ & ✓ & ✓ & ✓ (read) & ✗ \\
\bottomrule
\end{tabularx}
\end{table}

\textbf{Implementation:}
\begin{lstlisting}[style=pythonstyle]
# services/api-gateway/middleware/rbac.py
from functools import wraps
from fastapi import HTTPException, Depends

def require_role(allowed_roles: list):
    """Decorator to enforce role-based access"""
    def decorator(func):
        @wraps(func)
        async def wrapper(*args, current_user: dict = Depends(get_current_user), **kwargs):
            if current_user['role'] not in allowed_roles:
                raise HTTPException(
                    status_code=403,
                    detail=f"Access denied. Required roles: {allowed_roles}"
                )
            return await func(*args, current_user=current_user, **kwargs)
        return wrapper
    return decorator

# Usage example
@router.post("/api/v1/upload")
@require_role(['admin', 'data_manager'])
async def upload_documents(current_user: dict = Depends(get_current_user)):
    # Only admin and data_manager can upload
    pass
\end{lstlisting}

\subsection{Data Security}

\subsubsection{Encryption at Rest}

All sensitive data encrypted using industry-standard algorithms.

\begin{table}[h]
\centering
\caption{Encryption Standards by Data Type}
\begin{tabularx}{\textwidth}{lXl}
\toprule
\textbf{Data Type} & \textbf{Encryption Method} & \textbf{Key Length} \\
\midrule
PostgreSQL Database & AES-256-GCM & 256-bit \\
Document Content & AES-256-CTR & 256-bit \\
User Passwords & Argon2id (PHC winner) & N/A \\
API Keys & AES-256-GCM + Key Derivation & 256-bit \\
Session Data (Redis) & TLS 1.3 in transit & N/A \\
Elasticsearch Indices & XPack Security + TLS & 256-bit \\
Neo4j Graph Data & Enterprise Encryption & 256-bit \\
\bottomrule
\end{tabularx}
\end{table}

\textbf{Password Hashing Implementation:}
\begin{lstlisting}[style=pythonstyle]
# services/api-gateway/services/auth_service.py
from passlib.context import CryptContext

# Argon2id configuration (winner of Password Hashing Competition)
pwd_context = CryptContext(
    schemes=["argon2"],
    argon2__memory_cost=65536,  # 64 MB
    argon2__time_cost=3,        # 3 iterations
    argon2__parallelism=4       # 4 parallel threads
)

def hash_password(password: str) -> str:
    """Hash password using Argon2id"""
    return pwd_context.hash(password)

def verify_password(plain_password: str, hashed_password: str) -> bool:
    """Verify password against hash"""
    return pwd_context.verify(plain_password, hashed_password)
\end{lstlisting}

\subsubsection{Encryption in Transit}

\textbf{TLS/SSL Configuration:}
\begin{lstlisting}[style=pythonstyle]
# frontend/nginx-ssl.conf
server {
    listen 443 ssl http2;
    server_name yourdomain.com;
    
    # SSL certificates
    ssl_certificate /etc/ssl/certs/fullchain.pem;
    ssl_certificate_key /etc/ssl/private/privkey.pem;
    
    # Strong TLS configuration
    ssl_protocols TLSv1.3 TLSv1.2;
    ssl_ciphers 'ECDHE-ECDSA-AES256-GCM-SHA384:ECDHE-RSA-AES256-GCM-SHA384:ECDHE-ECDSA-CHACHA20-POLY1305';
    ssl_prefer_server_ciphers on;
    
    # Perfect Forward Secrecy
    ssl_ecdh_curve secp384r1;
    ssl_dhparam /etc/ssl/certs/dhparam.pem;
    
    # HSTS (HTTP Strict Transport Security)
    add_header Strict-Transport-Security "max-age=31536000; includeSubDomains; preload" always;
    
    # Security headers
    add_header X-Frame-Options "SAMEORIGIN" always;
    add_header X-Content-Type-Options "nosniff" always;
    add_header X-XSS-Protection "1; mode=block" always;
    add_header Referrer-Policy "strict-origin-when-cross-origin" always;
    add_header Content-Security-Policy "default-src 'self'; script-src 'self' 'unsafe-inline'; style-src 'self' 'unsafe-inline';" always;
}
\end{lstlisting}

\textbf{Inter-Service Communication:}
\begin{itemize}
    \item All microservice communication encrypted with mutual TLS (mTLS)
    \item Certificate-based authentication between services
    \item Automatic certificate rotation (Let's Encrypt)
    \item Service mesh ready (Istio/Linkerd compatible)
\end{itemize}

\subsection{Access Control Features}

\subsubsection{Multi-Factor Authentication (MFA)}

Optional 2FA using Time-based One-Time Passwords (TOTP).

\textbf{Implementation:}
\begin{lstlisting}[style=pythonstyle]
# services/api-gateway/services/mfa_service.py
import pyotp
import qrcode
from io import BytesIO

class MFAService:
    def generate_secret(self, user_email: str) -> dict:
        """Generate TOTP secret for user"""
        secret = pyotp.random_base32()
        totp = pyotp.TOTP(secret)
        
        # Generate provisioning URI for QR code
        provisioning_uri = totp.provisioning_uri(
            name=user_email,
            issuer_name="Data Ingestion Platform"
        )
        
        # Generate QR code
        qr = qrcode.QRCode(version=1, box_size=10, border=5)
        qr.add_data(provisioning_uri)
        qr.make(fit=True)
        
        img = qr.make_image(fill_color="black", back_color="white")
        buffer = BytesIO()
        img.save(buffer, format='PNG')
        qr_code_base64 = base64.b64encode(buffer.getvalue()).decode()
        
        return {
            "secret": secret,
            "qr_code": qr_code_base64,
            "provisioning_uri": provisioning_uri
        }
    
    def verify_totp(self, secret: str, token: str) -> bool:
        """Verify TOTP token"""
        totp = pyotp.TOTP(secret)
        return totp.verify(token, valid_window=1)  # Allow 30s drift
\end{lstlisting}

\textbf{MFA Features:}
\begin{itemize}
    \item Compatible with Google Authenticator, Authy, Microsoft Authenticator
    \item Backup codes generated during setup (10 one-time use codes)
    \item Account recovery via admin intervention
    \item Optional enforcement per role (mandatory for admins)
    \item Device trust (remember device for 30 days)
\end{itemize}

\begin{figure}[h]
\centering
\fbox{\textit{[SCREENSHOT PLACEHOLDER: MFA Setup - QR Code \& Verification]}}
\caption{Multi-Factor Authentication Setup Screen}
\label{fig:mfa-setup}
\end{figure}

\subsubsection{Session Management}

\begin{lstlisting}[style=pythonstyle]
# services/api-gateway/middleware/session.py
from redis import Redis
import hashlib

class SessionManager:
    def __init__(self, redis_client: Redis):
        self.redis = redis_client
        self.session_timeout = 1800  # 30 minutes
        self.max_sessions_per_user = 3  # Concurrent session limit
    
    def create_session(self, user_id: str, device_info: dict) -> str:
        """Create new session for user"""
        session_id = self._generate_session_id(user_id, device_info)
        
        # Check concurrent session limit
        active_sessions = self.get_active_sessions(user_id)
        if len(active_sessions) >= self.max_sessions_per_user:
            # Revoke oldest session
            self.revoke_session(active_sessions[0])
        
        # Store session data
        session_data = {
            "user_id": user_id,
            "created_at": datetime.utcnow().isoformat(),
            "device": device_info,
            "ip_address": device_info.get("ip"),
            "user_agent": device_info.get("user_agent")
        }
        
        self.redis.setex(
            f"session:{session_id}",
            self.session_timeout,
            json.dumps(session_data)
        )
        
        # Track active sessions for user
        self.redis.sadd(f"user_sessions:{user_id}", session_id)
        
        return session_id
    
    def validate_session(self, session_id: str) -> dict:
        """Validate and refresh session"""
        session_data = self.redis.get(f"session:{session_id}")
        if not session_data:
            raise HTTPException(status_code=401, detail="Session expired")
        
        # Refresh session timeout (sliding expiration)
        self.redis.expire(f"session:{session_id}", self.session_timeout)
        
        return json.loads(session_data)
    
    def revoke_session(self, session_id: str):
        """Revoke specific session"""
        session_data = self.redis.get(f"session:{session_id}")
        if session_data:
            user_id = json.loads(session_data)["user_id"]
            self.redis.srem(f"user_sessions:{user_id}", session_id)
        self.redis.delete(f"session:{session_id}")
\end{lstlisting}

\subsection{Audit \& Compliance}

\subsubsection{Comprehensive Audit Logging}

All user actions logged to PostgreSQL with immutable records.

\textbf{Audit Log Schema:}
\begin{lstlisting}[style=pythonstyle]
# services/api-gateway/models/audit_log.py
from sqlalchemy import Column, String, DateTime, JSON, Index

class AuditLog(Base):
    __tablename__ = "audit_logs"
    
    id = Column(String, primary_key=True)
    timestamp = Column(DateTime(timezone=True), nullable=False, index=True)
    user_id = Column(String, nullable=False, index=True)
    user_email = Column(String, nullable=False)
    action_type = Column(String, nullable=False, index=True)  # LOGIN, UPLOAD, CLASSIFY, etc.
    resource_type = Column(String, nullable=True)  # document, user, space
    resource_id = Column(String, nullable=True, index=True)
    action_details = Column(JSON, nullable=True)
    ip_address = Column(String, nullable=False)
    user_agent = Column(String, nullable=True)
    session_id = Column(String, nullable=False)
    status = Column(String, nullable=False)  # success, failure
    error_message = Column(String, nullable=True)
    
    __table_args__ = (
        Index('idx_audit_user_timestamp', 'user_id', 'timestamp'),
        Index('idx_audit_action_timestamp', 'action_type', 'timestamp'),
    )
\end{lstlisting}

\textbf{Logged Actions:}
\begin{itemize}
    \item \textbf{Authentication}: Login, logout, failed attempts, MFA verification
    \item \textbf{Document Operations}: Upload, download, delete, classify
    \item \textbf{User Management}: Create, update, delete users, role changes
    \item \textbf{Access Control}: Permission grants/revokes, ICSR approvals
    \item \textbf{Configuration}: System settings changes, API key generation
    \item \textbf{Data Access}: Document views, entity queries, API calls
\end{itemize}

\textbf{Compliance Features:}
\begin{table}[h]
\centering
\caption{Regulatory Compliance Standards}
\begin{tabularx}{\textwidth}{lX}
\toprule
\textbf{Standard} & \textbf{Compliance Status} \\
\midrule
\textbf{21 CFR Part 11} & ✓ Electronic signatures, audit trails, access control \\
\textbf{GDPR} & ✓ Data encryption, right to erasure, audit logs \\
\textbf{HIPAA} & ✓ PHI encryption, access controls, audit logs \\
\textbf{ISO 27001} & ✓ Information security management \\
\textbf{SOC 2 Type II} & ⏳ Planned Q2 2026 certification \\
\bottomrule
\end{tabularx}
\end{table}

\subsubsection{Data Privacy Features}

\begin{lstlisting}[style=pythonstyle]
# services/api-gateway/services/privacy_service.py
class PrivacyService:
    def anonymize_patient_data(self, text: str) -> str:
        """Redact PII from document text"""
        # Regex patterns for PII
        patterns = {
            'ssn': r'\b\d{3}-\d{2}-\d{4}\b',
            'email': r'\b[\w\.-]+@[\w\.-]+\.\w+\b',
            'phone': r'\b\d{3}[-.]?\d{3}[-.]?\d{4}\b',
            'mrn': r'\bMRN:\s*\d+\b',
            'dob': r'\b\d{1,2}/\d{1,2}/\d{2,4}\b'
        }
        
        for pii_type, pattern in patterns.items():
            text = re.sub(pattern, f'[{pii_type.upper()}_REDACTED]', text)
        
        return text
    
    def export_user_data(self, user_id: str) -> dict:
        """GDPR Article 15: Right to Access"""
        # Collect all user data from all systems
        user_data = {
            "profile": self.get_user_profile(user_id),
            "documents": self.get_user_documents(user_id),
            "audit_logs": self.get_user_audit_logs(user_id),
            "sessions": self.get_user_sessions(user_id)
        }
        return user_data
    
    def delete_user_data(self, user_id: str):
        """GDPR Article 17: Right to Erasure"""
        # Cascade delete user data (keeping audit trail)
        with transaction():
            # Anonymize user in audit logs (keep logs for compliance)
            self.anonymize_audit_logs(user_id)
            # Delete user documents
            self.delete_user_documents(user_id)
            # Delete user profile
            self.delete_user_profile(user_id)
            # Revoke all sessions
            self.revoke_all_sessions(user_id)
\end{lstlisting}

\subsection{API Security}

\subsubsection{Rate Limiting \& Throttling}

Multi-tier rate limiting to prevent abuse.

\begin{lstlisting}[style=pythonstyle]
# services/api-gateway/middleware/rate_limiter.py
from slowapi import Limiter
from slowapi.util import get_remote_address

limiter = Limiter(key_func=get_remote_address)

# Tier 1: Per-IP rate limiting (anonymous)
@limiter.limit("100/hour")
@router.get("/api/v1/public/search")
async def public_search(request: Request):
    pass

# Tier 2: Per-user rate limiting (authenticated)
@limiter.limit("1000/hour")
@router.get("/api/v1/documents")
async def get_documents(current_user: dict = Depends(get_current_user)):
    pass

# Tier 3: Cost-based limiting (for expensive operations)
@limiter.limit("10/minute")
@router.post("/api/v1/pv-classification/classify/batch")
async def classify_batch(current_user: dict = Depends(get_current_user)):
    pass

# Tier 4: Dynamic throttling based on system load
class AdaptiveRateLimiter:
    def __init__(self):
        self.redis = Redis()
    
    def get_current_limit(self, user_id: str) -> int:
        """Adjust rate limit based on system load"""
        cpu_usage = psutil.cpu_percent()
        
        if cpu_usage > 90:
            return 100  # Strict limit under high load
        elif cpu_usage > 70:
            return 500  # Moderate limit
        else:
            return 1000  # Normal limit
\end{lstlisting}

\subsubsection{Input Validation \& Sanitization}

\begin{lstlisting}[style=pythonstyle]
# services/api-gateway/validators/input_validator.py
from pydantic import BaseModel, validator, constr
import bleach

class DocumentUploadRequest(BaseModel):
    """Validate document upload request"""
    title: constr(min_length=1, max_length=500)
    space_id: constr(regex=r'^space_[a-zA-Z0-9]+$')
    
    @validator('title')
    def sanitize_title(cls, v):
        """Remove HTML/JS to prevent XSS"""
        return bleach.clean(v, tags=[], strip=True)

class UserCreateRequest(BaseModel):
    """Validate user creation"""
    email: EmailStr
    password: constr(min_length=12, max_length=128)
    role: str
    
    @validator('password')
    def validate_password_strength(cls, v):
        """Enforce strong password policy"""
        if not re.search(r'[A-Z]', v):
            raise ValueError('Password must contain uppercase')
        if not re.search(r'[a-z]', v):
            raise ValueError('Password must contain lowercase')
        if not re.search(r'[0-9]', v):
            raise ValueError('Password must contain digit')
        if not re.search(r'[!@#$%^&*]', v):
            raise ValueError('Password must contain special char')
        return v
    
    @validator('role')
    def validate_role(cls, v):
        """Ensure role is valid"""
        valid_roles = ['admin', 'data_manager', 'reviewer', 'analyst', 'guest']
        if v not in valid_roles:
            raise ValueError(f'Invalid role. Must be one of: {valid_roles}')
        return v
\end{lstlisting}

\section{Security Roadmap (2026-2027)}

\subsection{Q1 2026: Enhanced Authentication}

\textbf{Planned Features:}
\begin{enumerate}
    \item \textbf{Single Sign-On (SSO)}
    \begin{itemize}
        \item SAML 2.0 integration with corporate IdPs
        \item OAuth 2.0 / OpenID Connect support
        \item Microsoft Azure AD, Okta, Google Workspace
        \item Automatic user provisioning via SCIM
    \end{itemize}
    
    \item \textbf{Biometric Authentication}
    \begin{itemize}
        \item WebAuthn / FIDO2 support
        \item Fingerprint, Face ID for mobile devices
        \item Hardware security keys (YubiKey, Titan)
    \end{itemize}
    
    \item \textbf{Adaptive Authentication}
    \begin{itemize}
        \item Risk-based authentication (IP, device, behavior)
        \item Step-up authentication for sensitive actions
        \item Geo-location based access restrictions
    \end{itemize}
\end{enumerate}

\subsection{Q2 2026: Data Loss Prevention (DLP)}

\textbf{Planned Features:}
\begin{enumerate}
    \item \textbf{Content Inspection}
    \begin{itemize}
        \item Automatic PII detection in uploaded documents
        \item Configurable DLP policies (block, redact, alert)
        \item Integration with cloud DLP services (AWS Macie, Google DLP)
    \end{itemize}
    
    \item \textbf{Data Exfiltration Prevention}
    \begin{itemize}
        \item Document watermarking (visible \& invisible)
        \item Download restrictions per role
        \item Email gateway integration (prevent unauthorized sharing)
    \end{itemize}
    
    \item \textbf{Insider Threat Detection}
    \begin{itemize}
        \item ML-based anomaly detection (unusual access patterns)
        \item User behavior analytics (UEBA)
        \item Automated alerts to security team
    \end{itemize}
\end{enumerate}

\subsection{Q3 2026: Advanced Threat Protection}

\textbf{Planned Features:}
\begin{enumerate}
    \item \textbf{Intrusion Detection System (IDS)}
    \begin{itemize}
        \item Real-time network traffic analysis
        \item Signature-based + anomaly-based detection
        \item Integration with Suricata or Snort
    \end{itemize}
    
    \item \textbf{Malware Scanning}
    \begin{itemize}
        \item PDF malware analysis (macros, embedded scripts)
        \item Integration with VirusTotal API
        \item Sandbox execution for suspicious files
    \end{itemize}
    
    \item \textbf{Web Application Firewall (WAF)}
    \begin{itemize}
        \item OWASP Top 10 protection
        \item SQL injection, XSS, CSRF prevention
        \item Bot detection \& mitigation
    \end{itemize}
\end{enumerate}

\subsection{Q4 2026: Compliance \& Governance}

\textbf{Planned Features:}
\begin{enumerate}
    \item \textbf{SOC 2 Type II Certification}
    \begin{itemize}
        \item Full audit trail implementation
        \item Continuous monitoring dashboard
        \item Third-party audit preparation
    \end{itemize}
    
    \item \textbf{HIPAA Compliance}
    \begin{itemize}
        \item PHI encryption at rest \& in transit
        \item Business Associate Agreements (BAA)
        \item HIPAA audit logs \& reports
    \end{itemize}
    
    \item \textbf{Data Residency Controls}
    \begin{itemize}
        \item Multi-region deployment options
        \item Data sovereignty compliance (GDPR, CCPA)
        \item Configurable data retention policies
    \end{itemize}
\end{enumerate}

\subsection{Q1 2027: Zero Trust Architecture}

\textbf{Planned Features:}
\begin{enumerate}
    \item \textbf{Service Mesh Integration}
    \begin{itemize}
        \item Mutual TLS (mTLS) for all inter-service communication
        \item Service-to-service authentication
        \item Traffic encryption \& policy enforcement (Istio/Linkerd)
    \end{itemize}
    
    \item \textbf{Microsegmentation}
    \begin{itemize}
        \item Network isolation per microservice
        \item Fine-grained firewall rules
        \item Prevent lateral movement in case of breach
    \end{itemize}
    
    \item \textbf{Continuous Verification}
    \begin{itemize}
        \item Device posture checks before access
        \item Continuous authentication (not just login)
        \item Context-aware access decisions
    \end{itemize}
\end{enumerate}

\section{Security Best Practices}

\subsection{For Administrators}

\begin{infobox}
\textbf{Security Hardening Checklist:}
\begin{itemize}[noitemsep]
    \item Change all default passwords immediately
    \item Enable MFA for all admin accounts
    \item Rotate JWT secret keys quarterly
    \item Review audit logs weekly
    \item Keep all services updated (security patches)
    \item Restrict database access to localhost only
    \item Use strong TLS ciphers (disable TLS 1.0, 1.1)
    \item Enable firewall rules (only open required ports)
    \item Implement IP whitelisting for admin access
    \item Regularly backup encryption keys securely
\end{itemize}
\end{infobox}

\subsection{For End Users}

\begin{itemize}
    \item Use strong, unique passwords (16+ characters)
    \item Enable MFA if available
    \item Never share account credentials
    \item Log out from shared devices
    \item Report suspicious activity immediately
    \item Be cautious of phishing emails
    \item Verify document sources before uploading
    \item Don't bypass security warnings
\end{itemize}

\subsection{Security Incident Response}

\textbf{Incident Response Plan:}
\begin{enumerate}
    \item \textbf{Detection} - Automated alerts + user reports
    \item \textbf{Containment} - Isolate affected services
    \item \textbf{Eradication} - Remove threat, patch vulnerability
    \item \textbf{Recovery} - Restore services from backup
    \item \textbf{Post-Incident} - Root cause analysis, update procedures
\end{enumerate}

\textbf{Emergency Contacts:}
\begin{itemize}
    \item Security Team: security@auburn.edu
    \item On-Call Engineer: +1-555-SECURITY
    \item Incident Hotline: incidents@auburn.edu
\end{itemize}

\end{document}

